年次chronologyで最初に分けるべきなのはevent typeである。research paperの初出、organizationによるfirst-party publication、APIへのaccess開始、technology license、open-source / system releaseは互いに関連していても同一eventではない。本号は後年よく使われる単一の『release date』へ圧縮せず、それぞれが研究・distribution・利用権のどの境界を動かしたかで時系列を読む。


Q1には、scaleを経験則として扱う研究、large-scale open-domain dialogue、retrievalをpretrainingへ組み込む設計、frontier modelとlarge-model training systemが相次いで可視化された。ここではまだ一つのdominant product storyへ収束しておらず、『大きくする』『外部知識を使う』『会話を評価する』『訓練可能にする』という問題が別々の研究軸として並んでいる。


Q2は2020年の構造変化が最も密集する。dialogue system recipe、retrieval-conditioned generation、GPT-3のfew-shot framing、remote APIというaccess boundary、autoregressive image modeling、diffusion generation、distributed shardingが短い期間に重なった。重要なのは一つのmodelが全方向を統合したことではなく、model capabilityの周辺にaccess、retrieval、media、systemsという別layerが同時に立ち上がった点である。


Q3では、modelの『能力』以外の論点が前面へ出る。human preferenceをtraining loopへ入れるsteering、toxic degenerationのevaluation、first-party RAG boundary、GPT-3に関するlicense/access boundaryなどが、modelをどう使い、どう制御し、誰がどのsurfaceから利用できるかという問題を明確にした。scaleの議論がmodel本体だけで完結しなくなったことが見えやすい時期である。


Q4では、diffusion samplingのrefinement、score-based modelとSDEの統一的な整理、離散visual representationとTransformerを組み合わせるhigh-resolution synthesisなど、generative-media側の設計空間がさらに分化した。同時期のbias evaluationも含め、年末は『一つの巨大language modelへ収束した』のではなく、generation methodとevaluation layerが複数方向へ拡張した状態で2021年へ渡っていく。


このchronologyは四半期の構造を読むための本文であり、Evidenceで固定されていない正確な日付を現在の検索結果から補うものではない。最終的なDetailed Chronology tableでは、paper first-publication、first-party publication / release、API access、licenseを別eventとして保持し、確認済みのtimestampだけを記載する。後世の知名度で一つの日付を選び直さないことを優先する。
